
\emph{Epistemic Planning} (EP) refers to an automated planning setting where the agent reasons in the space of knowledge states and tries to find a plan to reach a desirable state from the current state. Its general form, the \emph{Multi-agent Epistemic Planning} (MEP) problem involves multiple agents who need to reason about both the state of the world and the information flow between agents.
In a MEP problem, multiple approaches have been developed recently with varying restrictions, such as considering only the concept of \emph{knowledge} while not allowing the idea of \emph{belief}, or not allowing for ``complex" modal operators such as those needed to handle \emph{dynamic common knowledge}. 
While the diversity of approaches has led to a deeper understanding of the problem space, the lack of a  standardized way to specify MEP problems independently of solution approaches has created difficulties in comparing performance of planners, identifying promising techniques, exploring new strategies like ensemble methods, and making it easy for new researchers to contribute to this research area.
To address the situation, we propose a unified way of specifying EP problems - the Epistemic Planning Domain Definition Language, E-PDDL. We show that E-PPDL can be supported by leading MEP planners and provide corresponding parser code that translates EP problems specified in E-PDDL into (M)EP problems that can be handled by several planners. 
This work is also useful in building more general epistemic planning environments where we envision a meta-cognitive module that takes a planning problem in E-PDDL, identifies and assesses some of its features, and autonomously decides which planner is the best one to solve it.